\documentclass[t]{beamer}
\usepackage[utf8]{inputenc}
\usepackage[francais]{babel}
\usepackage{tikz,pgf}
\usepackage{wrapfig}
\usepackage{times}
\usepackage[absolute,overlay]{textpos}
\usepackage{url}
% \usepackage[texcoord,grid,gridunit=mm,gridcolor=red!10,subgridcolor=green!10]{eso-pic}

\definecolor{bleu}{rgb}{0.282, 0.439, 0.643}

\usecolortheme[named=bleu]{structure}
\useinnertheme[shadow]{rounded}

\hypersetup{pdfpagemode=FullScreen}

\setbeamersize{text margin left=0.6cm}
\setbeamersize{text margin right=0.4cm}
\setbeamersize{text margin left=.5cm, text margin right=.5cm}
\setbeamertemplate{navigation symbols}{}

\AtBeginSection[]{
    \begin{frame}
  	 \frametitle{\vspace*{-1cm} \begin{flushright} \ \ \vskip-0.25cm \end{flushright} }
        \vspace*{-0.75cm}
	    \begin{center}
            \centering
            \tableofcontents[currentsection]
	    \end{center}
    \end{frame}
}

\begin{document}
    \begin{frame}
        \pgfputat{\pgfxy(-0.505,0.685)}{\pgfbox[left,top]{\pgfimage[width=\paperwidth,height=\paperheight ]{./assets/UTT_garde_empty.pdf}}}
        \vspace*{2.5cm}

        \begin{center}
            \textcolor{white}{
                {\large Soutenance de TITS} \\
                {\footnotesize   Printemps 2015}\\
            \vspace*{0.5cm}
                {\footnotesize \textbf{Développement d'un robot autonome}}\\
            \vspace*{0.25cm}
                {\footnotesize Mousset Axel\\ Labate Aurélien}\\
          }
        \end{center}
    \end{frame}

\setbeamertemplate{background canvas}{
    \includegraphics[width=\paperwidth,height=\paperheight]{./assets/UTT_fond.pdf}}
    \addtobeamertemplate{footline}{
    \vspace{-0.5cm}
    \begin{flushright}
    \insertframenumber/\inserttotalframenumber
    \end{flushright}
}


\section{Introduction}
    \begin{frame}[c]
        \begin{center}
            \frametitle{\vspace*{-1.0cm} \begin{flushright} Introduction \ \ \vskip-0.25cm \end{flushright} }
            \vspace*{-0.75cm}
            \begin{itemize}
                \item Développement d'un robot autonome pour la coupe de France
                \item Classement: 54èmes sur plus de 180 équipes
                \item Casse mécanique du robot principal !
                \item Objectif: proposer des bases théoriques, algorithmiques et logicielles solides pour les années suivantes
            \end{itemize}
        \end{center}
    \end{frame}

    \begin{frame}[c]
        \frametitle{\vspace*{-1.0cm} \begin{flushright} Introduction \ \ \vskip-0.25cm \end{flushright} }
        La robotique est une discipline complexe, elle allie:
        \begin{center}
            \begin{itemize}
                \item Informatique
                \item Électronique
                \item Mécanique
                \item Sciences physiques, algorithmique..
            \end{itemize}
        \end{center}
        Dans cette présentation, on parlera d'\textbf{informatique} et de \textbf{réseau}.
    \end{frame}


\section{Architecture modulaire}
    \begin{frame}[c]
        \frametitle{\vspace*{-1.0cm} \begin{flushright} Architecture modulaire \ \ \vskip-0.25cm \end{flushright} }
        \vspace*{-0.75cm}
    \end{frame}


\section{Technologies de communication}
        \begin{frame}[c]
            \frametitle{\vspace*{-1.0cm} \begin{flushright} Technologies de communication \ \ \vskip-0.25cm \end{flushright} }
            \vspace*{0.5cm}
             \begin{textblock*}{1cm}(4cm,3cm)
                \textbf{UART}
            \end{textblock*}
            \begin{center}
                \begin{itemize}
                    \item Liaison full-duplex
                    \item Niveau logique: TTL
                    \item Connexion asynchrone
                    \item CRC basique: bit de parité
                    \item Débit (baudrate) variable, maximum 3,6GBauds théoriques
                    \item 2 fils par pair de périphérique
                \end{itemize}
            \end{center}
        \end{frame}

        \begin{frame}[c]
            \frametitle{\vspace*{-1.0cm} \begin{flushright} Technologies de communication \ \ \vskip-0.25cm \end{flushright} }
            \vspace*{0.5cm}
             \begin{textblock*}{1cm}(4cm,3cm)
                \textbf{SPI}
            \end{textblock*}
            \begin{center}
                \begin{itemize}
                    \item Liaison full-duplex
                    \item Bus de donnée, relation master/slave
                    \item Connexion synchrone: le master donne la clock
                    \item Bus de donnée de deux fil, et un fil d'addressage par esclave
                    \item Débit élevé pour un bus, jusqu'à 20Mb/s
                    \item Protocle simple, pas d'adressage/gestion d'erreur
                \end{itemize}
            \end{center}
        \end{frame}

        \begin{frame}[c]
            \frametitle{\vspace*{-1.0cm} \begin{flushright} Technologies de communication \ \ \vskip-0.25cm \end{flushright} }
            \vspace*{0.5cm}
             \begin{textblock*}{1cm}(4cm,3cm)
                \textbf{CAN}
            \end{textblock*}
            \begin{center}
                \begin{itemize}
                    \item Liaison half-duplex
                    \item Bus de donnée, relation multi-master
                    \item Connexion asynchrone
                    \item Detection et retour sur erreur physique
                    \item Débit moyen pour un bus, jusqu'à 1Mb/s
                    \item Fonctionnement "événementiel"
                    \item Jusqu'à 30 périphériques, système d'adressage
                    \item Très résistant aux perturbations
                \end{itemize}
            \end{center}
        \end{frame}

        \begin{frame}[c]
            \frametitle{\vspace*{-1.0cm} \begin{flushright} Technologies de communication \ \ \vskip-0.25cm \end{flushright} }
            \vspace*{0.5cm}
             \begin{textblock*}{1cm}(4cm,3cm)
                \textbf{I2C}
            \end{textblock*}
            \begin{center}
                \begin{itemize}
                    \item Liaison half-duplex
                    \item Bus de donnée, relation master/slave ou multi-master
                    \item Connexion synchrone: un fil commun de clock
                    \item Débit faible, jusqu'à 400Kb/s
                    \item Jusqu'à une centaine de périphériques, système d'adressage
                    \item Un bus de donnée sur un fil
                    \item Aquittement des données
                \end{itemize}
            \end{center}
        \end{frame}

        \begin{frame}[c]
            \frametitle{\vspace*{-1.0cm} \begin{flushright} Technologies de communication \ \ \vskip-0.25cm \end{flushright} }
            \vspace*{0.5cm}
             \begin{textblock*}{1cm}(4cm,3cm)
            \end{textblock*}
            Technologie choisie: I2C.\\
            Problèmes:
            \begin{center}
                \begin{itemize}
                    \item Intégrité des données
                    \item Architecture master/slave
                    \item Librairies Javascript peu abouties
                \end{itemize}
            \end{center}


        \begin{frame}[c]
            \frametitle{\vspace*{-1.0cm} \begin{flushright} Technologies de communication \ \ \vskip-0.25cm \end{flushright} }
            \vspace*{0.5cm}
            \begin{textblock*}{1cm}(4cm,3cm)
             Construction d'une surcouche du protocole
            \end{textblock*}
            Le principal but est d'avoir un retour sur erreur.
            \begin{center}
                \begin{itemize}
                    \item Taille du packet annoncée
                    \item Numéro de packet "Fonction"
                    \item Données "Arguments"
                    \item CRC: XORSum
                \end{itemize}
                INSÉRER UNE TABLEA-TRAME ICI
            \end{center}
        \end{frame}

\section{Conclusion}
    \begin{frame}[c]
        \frametitle{\vspace*{-1.0cm} \begin{flushright} Conclusion \ \ \vskip-0.25cm \end{flushright} }
        \vspace*{-0.75cm}
    \end{frame}

\section{Démonstration}
    \begin{frame}[c]
        \frametitle{\vspace*{-1.0cm} \begin{flushright} Démonstration \ \ \vskip-0.25cm \end{flushright} }
        \vspace*{-0.75cm}
        \begin{center}
            \href{./assets/demo.mp4}{Vidéo}
        \end{center}
    \end{frame}



\end{document}